\section{Formulación matemática}
Sea $\mathcal I=\{1,\ldots,15\}$ el conjunto de embarcaciones y $\mathcal J$ el conjunto de zonas candidatas. La variable de decisión es:
\[
x_{ij}=\begin{cases}1,&\text{si el barco }i\text{ se asigna a la zona }j,\\0,&\text{en otro caso.}\end{cases}
\]
Para cada zona $j$ se conoce $P_j\in[0,1]$ y para cada par puerto-zona se calcula una distancia $D_{ij}$. El combustible homogéneo se expresa como:
\[F_{ij}=cD_{ij},\]
donde $c$ es común a toda la flota.
La función objetivo básica es:
\[
\max_x\sum_{i\in\mathcal I}\sum_{j\in\mathcal J}\left(\alpha P_j-\lambda \widetilde D_{ij}\right)x_{ij},
\]
sujeta a:
\[\sum_j x_{ij}=1\quad\forall i,\]
\[\sum_i x_{ij}\le K\quad\forall j,\]
\[x_{ij}\le L_j,\qquad x_{ij}\in\{0,1\}.\]
Para el componente multiagente se emplea una recompensa:
\[r_i=\alpha P_{a_i}-\lambda \widetilde D_{i,a_i}-\gamma C_{a_i}-\eta\,\mathrm{Fairness}.\]
El término $C_{a_i}$ penaliza congestión y el componente de equidad evita que el esfuerzo de navegación se concentre de forma desbalanceada entre agentes.